% paper39-generated-contracts.tex — Automatic Contract Generation for Python Libraries
% Paper 39 of the Judgment Geometry series.
% Compile: pdflatex -interaction=nonstopmode paper39-generated-contracts.tex
\documentclass[11pt]{article}
\usepackage{lmodern}
% jugeo-common.tex — shared preamble for all Judgment Geometry papers
% Include via: % jugeo-common.tex — shared preamble for all Judgment Geometry papers
% Include via: % jugeo-common.tex — shared preamble for all Judgment Geometry papers
% Include via: \input{jugeo-common}

% ── Packages ────────────────────────────────────────────────────────────
\usepackage{amsmath,amssymb,amsthm,mathtools}
\usepackage{tikz-cd}
\usepackage{listings}
\usepackage{xcolor}
\usepackage{xspace}
\usepackage{booktabs}
\usepackage{multirow}
\usepackage{enumitem}
\usepackage[expansion=false]{microtype}
\usepackage{hyperref}
\usepackage{cleveref}
% \usepackage{thmtools}  % disabled: not available in this TeX installation
\usepackage{float}
\usepackage{caption}
\usepackage{subcaption}
\usepackage{tabularx}
\usepackage{stmaryrd}       % \llbracket, \rrbracket
\usepackage{bbm}            % \mathbbm{1}

% ── Page geometry ───────────────────────────────────────────────────────
\usepackage[margin=1in]{geometry}

% ── Theorem environments ────────────────────────────────────────────────
\theoremstyle{plain}
\newtheorem{theorem}{Theorem}[section]
\newtheorem{lemma}[theorem]{Lemma}
\newtheorem{proposition}[theorem]{Proposition}
\newtheorem{corollary}[theorem]{Corollary}

\theoremstyle{definition}
\newtheorem{definition}[theorem]{Definition}
\newtheorem{example}[theorem]{Example}
\newtheorem{construction}[theorem]{Construction}

\theoremstyle{remark}
\newtheorem{remark}[theorem]{Remark}
\newtheorem{notation}[theorem]{Notation}

% ── Python listings style ──────────────────────────────────────────────
\definecolor{jugeoBlue}{HTML}{2563EB}
\definecolor{jugeoGreen}{HTML}{059669}
\definecolor{jugeoRed}{HTML}{DC2626}
\definecolor{jugeoPurple}{HTML}{7C3AED}
\definecolor{jugeoGray}{HTML}{6B7280}
\definecolor{jugeoBg}{HTML}{F8FAFC}

\lstdefinestyle{jugeo-python}{
  language=Python,
  basicstyle=\ttfamily\small,
  keywordstyle=\color{jugeoBlue}\bfseries,
  stringstyle=\color{jugeoGreen},
  commentstyle=\color{jugeoGray}\itshape,
  numberstyle=\tiny\color{jugeoGray},
  backgroundcolor=\color{jugeoBg},
  numbers=left,
  numbersep=6pt,
  frame=single,
  framerule=0.4pt,
  rulecolor=\color{jugeoGray!40},
  breaklines=true,
  breakatwhitespace=true,
  showstringspaces=false,
  tabsize=4,
  xleftmargin=1.5em,
  framexleftmargin=1.5em,
  morekeywords={dataclass,frozen,field,Enum,IntEnum,auto,Any,
                Mapping,Sequence,Optional,match,case},
  literate={→}{{$\to$}}1 {←}{{$\leftarrow$}}1
           {≤}{{$\leq$}}1 {≥}{{$\geq$}}1 {≠}{{$\neq$}}1
           {∈}{{$\in$}}1 {∀}{{$\forall$}}1 {∃}{{$\exists$}}1
           {⊕}{{$\oplus$}}1 {⊖}{{$\ominus$}}1,
  escapeinside={(*@}{@*)},
}
\lstset{style=jugeo-python}

\lstdefinestyle{jugeo-output}{
  basicstyle=\ttfamily\small,
  backgroundcolor=\color{jugeoBg},
  frame=single,
  framerule=0.4pt,
  rulecolor=\color{jugeoGray!40},
  breaklines=true,
  showstringspaces=false,
  xleftmargin=1.5em,
  framexleftmargin=0.5em,
  numbers=none,
}

% ── JuGeo-specific macros ──────────────────────────────────────────────

% Judgment 8-tuple
\newcommand{\J}{\mathcal{J}}
\newcommand{\Jud}[8]{(#1,\,#2,\,#3,\,#4,\,#5,\,#6,\,#7,\,#8)}
\newcommand{\JudShort}[1]{\J_{#1}}

% Components
\newcommand{\coord}{\mathsf{c}}            % coordinate
\newcommand{\prop}{\varphi}                 % proposition
\newcommand{\carrier}{\mathcal{A}}          % carrier type
\newcommand{\evid}{\mathcal{E}}             % evidence bundle
\newcommand{\oblig}{\mathcal{O}}            % obligations
\newcommand{\obstr}{\mathcal{B}}            % obstructions
\newcommand{\trust}{\mathcal{T}}            % trust annotation
\newcommand{\prov}{\Pi}                     % provenance

% Trust algebra
\newcommand{\Talg}{\mathfrak{T}}            % trust ordered algebra
\newcommand{\Eadm}{\mathcal{E}_{\mathrm{adm}}}
\newcommand{\tleq}{\preceq}                 % trust ordering
\newcommand{\tjoin}{\oplus}                 % conservative join
\newcommand{\tatten}{\ominus}               % attenuation
\newcommand{\tpromote}{\uparrow_\pi}        % promotion
\newcommand{\tdemote}{\downarrow_\chi}       % demotion

% Trust tiers
\newcommand{\Tcontra}{\ensuremath{\mathsf{CONTRADICTED}}}
\newcommand{\Tunver}{\ensuremath{\mathsf{UNVERIFIED}}}
\newcommand{\Tcopilot}{\ensuremath{\mathsf{COPILOT}}}
\newcommand{\Toracle}{\ensuremath{\mathsf{ORACLE}}}
\newcommand{\Truntime}{\ensuremath{\mathsf{RUNTIME}}}
\newcommand{\Tsolver}{\ensuremath{\mathsf{SOLVER}}}
\newcommand{\Tproof}{\ensuremath{\mathsf{PROOF}}}

% Site and geometry
\newcommand{\Site}{\mathcal{S}}             % semantic site
\newcommand{\Cov}{\mathrm{Cov}}            % covering family
\newcommand{\Ob}{\mathrm{Ob}}              % objects
\newcommand{\Mor}{\mathrm{Mor}}            % morphisms
\newcommand{\res}{\mathrm{res}}            % restriction
\newcommand{\Cech}{\check{C}}              % Čech complex
\newcommand{\Hcoh}[1]{H^{#1}}             % cohomology group

% Descent
\newcommand{\descent}{\mathrm{desc}}
\newcommand{\glue}{\mathrm{glue}}
\newcommand{\overlap}[2]{#1 \cap #2}

% Semantic moves
\newcommand{\VERIFY}{\textsc{Verify}}
\newcommand{\CONSTRUCT}{\textsc{Construct}}
\newcommand{\NORMALIZE}{\textsc{Normalize}}
\newcommand{\GLUE}{\textsc{Glue}}
\newcommand{\RESTRICT}{\textsc{Restrict}}
\newcommand{\TRANSPORT}{\textsc{Transport}}
\newcommand{\REPAIR}{\textsc{Repair}}
\newcommand{\PROMOTE}{\textsc{Promote}}
\newcommand{\CHALLENGE}{\textsc{Challenge}}
\newcommand{\ESCALATE}{\textsc{Escalate}}

% Fragments
\newcommand{\QFLIA}{\texttt{QF\_LIA}}
\newcommand{\QFLRA}{\texttt{QF\_LRA}}
\newcommand{\QFBV}{\texttt{QF\_BV}}
\newcommand{\QFUF}{\texttt{QF\_UF}}

% Misc
\newcommand{\Python}{\textsc{Python}}
\newcommand{\JuGeo}{\textsc{JuGeo}}
\newcommand{\LEAN}{\textsc{Lean}\,4}
\newcommand{\Fstar}{F$^{\star}$}
\newcommand{\Coq}{\textsc{Coq}}
\newcommand{\Dafny}{\textsc{Dafny}}

% Common shortcuts
\newcommand{\ie}{i.e.\@\xspace}
\newcommand{\eg}{e.g.\@\xspace}
\newcommand{\cf}{cf.\@\xspace}
\newcommand{\wrt}{w.r.t.\@\xspace}

% Evaluation shortcuts
\newcommand{\numcases}{300}
\newcommand{\numspec}{100}
\newcommand{\numeq}{100}
\newcommand{\numbug}{100}
\newcommand{\medtime}{6.5\,ms}
\newcommand{\totaltime}{1.949\,s}


% ── Packages ────────────────────────────────────────────────────────────
\usepackage{amsmath,amssymb,amsthm,mathtools}
\usepackage{tikz-cd}
\usepackage{listings}
\usepackage{xcolor}
\usepackage{xspace}
\usepackage{booktabs}
\usepackage{multirow}
\usepackage{enumitem}
\usepackage[expansion=false]{microtype}
\usepackage{hyperref}
\usepackage{cleveref}
% \usepackage{thmtools}  % disabled: not available in this TeX installation
\usepackage{float}
\usepackage{caption}
\usepackage{subcaption}
\usepackage{tabularx}
\usepackage{stmaryrd}       % \llbracket, \rrbracket
\usepackage{bbm}            % \mathbbm{1}

% ── Page geometry ───────────────────────────────────────────────────────
\usepackage[margin=1in]{geometry}

% ── Theorem environments ────────────────────────────────────────────────
\theoremstyle{plain}
\newtheorem{theorem}{Theorem}[section]
\newtheorem{lemma}[theorem]{Lemma}
\newtheorem{proposition}[theorem]{Proposition}
\newtheorem{corollary}[theorem]{Corollary}

\theoremstyle{definition}
\newtheorem{definition}[theorem]{Definition}
\newtheorem{example}[theorem]{Example}
\newtheorem{construction}[theorem]{Construction}

\theoremstyle{remark}
\newtheorem{remark}[theorem]{Remark}
\newtheorem{notation}[theorem]{Notation}

% ── Python listings style ──────────────────────────────────────────────
\definecolor{jugeoBlue}{HTML}{2563EB}
\definecolor{jugeoGreen}{HTML}{059669}
\definecolor{jugeoRed}{HTML}{DC2626}
\definecolor{jugeoPurple}{HTML}{7C3AED}
\definecolor{jugeoGray}{HTML}{6B7280}
\definecolor{jugeoBg}{HTML}{F8FAFC}

\lstdefinestyle{jugeo-python}{
  language=Python,
  basicstyle=\ttfamily\small,
  keywordstyle=\color{jugeoBlue}\bfseries,
  stringstyle=\color{jugeoGreen},
  commentstyle=\color{jugeoGray}\itshape,
  numberstyle=\tiny\color{jugeoGray},
  backgroundcolor=\color{jugeoBg},
  numbers=left,
  numbersep=6pt,
  frame=single,
  framerule=0.4pt,
  rulecolor=\color{jugeoGray!40},
  breaklines=true,
  breakatwhitespace=true,
  showstringspaces=false,
  tabsize=4,
  xleftmargin=1.5em,
  framexleftmargin=1.5em,
  morekeywords={dataclass,frozen,field,Enum,IntEnum,auto,Any,
                Mapping,Sequence,Optional,match,case},
  literate={→}{{$\to$}}1 {←}{{$\leftarrow$}}1
           {≤}{{$\leq$}}1 {≥}{{$\geq$}}1 {≠}{{$\neq$}}1
           {∈}{{$\in$}}1 {∀}{{$\forall$}}1 {∃}{{$\exists$}}1
           {⊕}{{$\oplus$}}1 {⊖}{{$\ominus$}}1,
  escapeinside={(*@}{@*)},
}
\lstset{style=jugeo-python}

\lstdefinestyle{jugeo-output}{
  basicstyle=\ttfamily\small,
  backgroundcolor=\color{jugeoBg},
  frame=single,
  framerule=0.4pt,
  rulecolor=\color{jugeoGray!40},
  breaklines=true,
  showstringspaces=false,
  xleftmargin=1.5em,
  framexleftmargin=0.5em,
  numbers=none,
}

% ── JuGeo-specific macros ──────────────────────────────────────────────

% Judgment 8-tuple
\newcommand{\J}{\mathcal{J}}
\newcommand{\Jud}[8]{(#1,\,#2,\,#3,\,#4,\,#5,\,#6,\,#7,\,#8)}
\newcommand{\JudShort}[1]{\J_{#1}}

% Components
\newcommand{\coord}{\mathsf{c}}            % coordinate
\newcommand{\prop}{\varphi}                 % proposition
\newcommand{\carrier}{\mathcal{A}}          % carrier type
\newcommand{\evid}{\mathcal{E}}             % evidence bundle
\newcommand{\oblig}{\mathcal{O}}            % obligations
\newcommand{\obstr}{\mathcal{B}}            % obstructions
\newcommand{\trust}{\mathcal{T}}            % trust annotation
\newcommand{\prov}{\Pi}                     % provenance

% Trust algebra
\newcommand{\Talg}{\mathfrak{T}}            % trust ordered algebra
\newcommand{\Eadm}{\mathcal{E}_{\mathrm{adm}}}
\newcommand{\tleq}{\preceq}                 % trust ordering
\newcommand{\tjoin}{\oplus}                 % conservative join
\newcommand{\tatten}{\ominus}               % attenuation
\newcommand{\tpromote}{\uparrow_\pi}        % promotion
\newcommand{\tdemote}{\downarrow_\chi}       % demotion

% Trust tiers
\newcommand{\Tcontra}{\ensuremath{\mathsf{CONTRADICTED}}}
\newcommand{\Tunver}{\ensuremath{\mathsf{UNVERIFIED}}}
\newcommand{\Tcopilot}{\ensuremath{\mathsf{COPILOT}}}
\newcommand{\Toracle}{\ensuremath{\mathsf{ORACLE}}}
\newcommand{\Truntime}{\ensuremath{\mathsf{RUNTIME}}}
\newcommand{\Tsolver}{\ensuremath{\mathsf{SOLVER}}}
\newcommand{\Tproof}{\ensuremath{\mathsf{PROOF}}}

% Site and geometry
\newcommand{\Site}{\mathcal{S}}             % semantic site
\newcommand{\Cov}{\mathrm{Cov}}            % covering family
\newcommand{\Ob}{\mathrm{Ob}}              % objects
\newcommand{\Mor}{\mathrm{Mor}}            % morphisms
\newcommand{\res}{\mathrm{res}}            % restriction
\newcommand{\Cech}{\check{C}}              % Čech complex
\newcommand{\Hcoh}[1]{H^{#1}}             % cohomology group

% Descent
\newcommand{\descent}{\mathrm{desc}}
\newcommand{\glue}{\mathrm{glue}}
\newcommand{\overlap}[2]{#1 \cap #2}

% Semantic moves
\newcommand{\VERIFY}{\textsc{Verify}}
\newcommand{\CONSTRUCT}{\textsc{Construct}}
\newcommand{\NORMALIZE}{\textsc{Normalize}}
\newcommand{\GLUE}{\textsc{Glue}}
\newcommand{\RESTRICT}{\textsc{Restrict}}
\newcommand{\TRANSPORT}{\textsc{Transport}}
\newcommand{\REPAIR}{\textsc{Repair}}
\newcommand{\PROMOTE}{\textsc{Promote}}
\newcommand{\CHALLENGE}{\textsc{Challenge}}
\newcommand{\ESCALATE}{\textsc{Escalate}}

% Fragments
\newcommand{\QFLIA}{\texttt{QF\_LIA}}
\newcommand{\QFLRA}{\texttt{QF\_LRA}}
\newcommand{\QFBV}{\texttt{QF\_BV}}
\newcommand{\QFUF}{\texttt{QF\_UF}}

% Misc
\newcommand{\Python}{\textsc{Python}}
\newcommand{\JuGeo}{\textsc{JuGeo}}
\newcommand{\LEAN}{\textsc{Lean}\,4}
\newcommand{\Fstar}{F$^{\star}$}
\newcommand{\Coq}{\textsc{Coq}}
\newcommand{\Dafny}{\textsc{Dafny}}

% Common shortcuts
\newcommand{\ie}{i.e.\@\xspace}
\newcommand{\eg}{e.g.\@\xspace}
\newcommand{\cf}{cf.\@\xspace}
\newcommand{\wrt}{w.r.t.\@\xspace}

% Evaluation shortcuts
\newcommand{\numcases}{300}
\newcommand{\numspec}{100}
\newcommand{\numeq}{100}
\newcommand{\numbug}{100}
\newcommand{\medtime}{6.5\,ms}
\newcommand{\totaltime}{1.949\,s}


% ── Packages ────────────────────────────────────────────────────────────
\usepackage{amsmath,amssymb,amsthm,mathtools}
\usepackage{tikz-cd}
\usepackage{listings}
\usepackage{xcolor}
\usepackage{xspace}
\usepackage{booktabs}
\usepackage{multirow}
\usepackage{enumitem}
\usepackage[expansion=false]{microtype}
\usepackage{hyperref}
\usepackage{cleveref}
% \usepackage{thmtools}  % disabled: not available in this TeX installation
\usepackage{float}
\usepackage{caption}
\usepackage{subcaption}
\usepackage{tabularx}
\usepackage{stmaryrd}       % \llbracket, \rrbracket
\usepackage{bbm}            % \mathbbm{1}

% ── Page geometry ───────────────────────────────────────────────────────
\usepackage[margin=1in]{geometry}

% ── Theorem environments ────────────────────────────────────────────────
\theoremstyle{plain}
\newtheorem{theorem}{Theorem}[section]
\newtheorem{lemma}[theorem]{Lemma}
\newtheorem{proposition}[theorem]{Proposition}
\newtheorem{corollary}[theorem]{Corollary}

\theoremstyle{definition}
\newtheorem{definition}[theorem]{Definition}
\newtheorem{example}[theorem]{Example}
\newtheorem{construction}[theorem]{Construction}

\theoremstyle{remark}
\newtheorem{remark}[theorem]{Remark}
\newtheorem{notation}[theorem]{Notation}

% ── Python listings style ──────────────────────────────────────────────
\definecolor{jugeoBlue}{HTML}{2563EB}
\definecolor{jugeoGreen}{HTML}{059669}
\definecolor{jugeoRed}{HTML}{DC2626}
\definecolor{jugeoPurple}{HTML}{7C3AED}
\definecolor{jugeoGray}{HTML}{6B7280}
\definecolor{jugeoBg}{HTML}{F8FAFC}

\lstdefinestyle{jugeo-python}{
  language=Python,
  basicstyle=\ttfamily\small,
  keywordstyle=\color{jugeoBlue}\bfseries,
  stringstyle=\color{jugeoGreen},
  commentstyle=\color{jugeoGray}\itshape,
  numberstyle=\tiny\color{jugeoGray},
  backgroundcolor=\color{jugeoBg},
  numbers=left,
  numbersep=6pt,
  frame=single,
  framerule=0.4pt,
  rulecolor=\color{jugeoGray!40},
  breaklines=true,
  breakatwhitespace=true,
  showstringspaces=false,
  tabsize=4,
  xleftmargin=1.5em,
  framexleftmargin=1.5em,
  morekeywords={dataclass,frozen,field,Enum,IntEnum,auto,Any,
                Mapping,Sequence,Optional,match,case},
  literate={→}{{$\to$}}1 {←}{{$\leftarrow$}}1
           {≤}{{$\leq$}}1 {≥}{{$\geq$}}1 {≠}{{$\neq$}}1
           {∈}{{$\in$}}1 {∀}{{$\forall$}}1 {∃}{{$\exists$}}1
           {⊕}{{$\oplus$}}1 {⊖}{{$\ominus$}}1,
  escapeinside={(*@}{@*)},
}
\lstset{style=jugeo-python}

\lstdefinestyle{jugeo-output}{
  basicstyle=\ttfamily\small,
  backgroundcolor=\color{jugeoBg},
  frame=single,
  framerule=0.4pt,
  rulecolor=\color{jugeoGray!40},
  breaklines=true,
  showstringspaces=false,
  xleftmargin=1.5em,
  framexleftmargin=0.5em,
  numbers=none,
}

% ── JuGeo-specific macros ──────────────────────────────────────────────

% Judgment 8-tuple
\newcommand{\J}{\mathcal{J}}
\newcommand{\Jud}[8]{(#1,\,#2,\,#3,\,#4,\,#5,\,#6,\,#7,\,#8)}
\newcommand{\JudShort}[1]{\J_{#1}}

% Components
\newcommand{\coord}{\mathsf{c}}            % coordinate
\newcommand{\prop}{\varphi}                 % proposition
\newcommand{\carrier}{\mathcal{A}}          % carrier type
\newcommand{\evid}{\mathcal{E}}             % evidence bundle
\newcommand{\oblig}{\mathcal{O}}            % obligations
\newcommand{\obstr}{\mathcal{B}}            % obstructions
\newcommand{\trust}{\mathcal{T}}            % trust annotation
\newcommand{\prov}{\Pi}                     % provenance

% Trust algebra
\newcommand{\Talg}{\mathfrak{T}}            % trust ordered algebra
\newcommand{\Eadm}{\mathcal{E}_{\mathrm{adm}}}
\newcommand{\tleq}{\preceq}                 % trust ordering
\newcommand{\tjoin}{\oplus}                 % conservative join
\newcommand{\tatten}{\ominus}               % attenuation
\newcommand{\tpromote}{\uparrow_\pi}        % promotion
\newcommand{\tdemote}{\downarrow_\chi}       % demotion

% Trust tiers
\newcommand{\Tcontra}{\ensuremath{\mathsf{CONTRADICTED}}}
\newcommand{\Tunver}{\ensuremath{\mathsf{UNVERIFIED}}}
\newcommand{\Tcopilot}{\ensuremath{\mathsf{COPILOT}}}
\newcommand{\Toracle}{\ensuremath{\mathsf{ORACLE}}}
\newcommand{\Truntime}{\ensuremath{\mathsf{RUNTIME}}}
\newcommand{\Tsolver}{\ensuremath{\mathsf{SOLVER}}}
\newcommand{\Tproof}{\ensuremath{\mathsf{PROOF}}}

% Site and geometry
\newcommand{\Site}{\mathcal{S}}             % semantic site
\newcommand{\Cov}{\mathrm{Cov}}            % covering family
\newcommand{\Ob}{\mathrm{Ob}}              % objects
\newcommand{\Mor}{\mathrm{Mor}}            % morphisms
\newcommand{\res}{\mathrm{res}}            % restriction
\newcommand{\Cech}{\check{C}}              % Čech complex
\newcommand{\Hcoh}[1]{H^{#1}}             % cohomology group

% Descent
\newcommand{\descent}{\mathrm{desc}}
\newcommand{\glue}{\mathrm{glue}}
\newcommand{\overlap}[2]{#1 \cap #2}

% Semantic moves
\newcommand{\VERIFY}{\textsc{Verify}}
\newcommand{\CONSTRUCT}{\textsc{Construct}}
\newcommand{\NORMALIZE}{\textsc{Normalize}}
\newcommand{\GLUE}{\textsc{Glue}}
\newcommand{\RESTRICT}{\textsc{Restrict}}
\newcommand{\TRANSPORT}{\textsc{Transport}}
\newcommand{\REPAIR}{\textsc{Repair}}
\newcommand{\PROMOTE}{\textsc{Promote}}
\newcommand{\CHALLENGE}{\textsc{Challenge}}
\newcommand{\ESCALATE}{\textsc{Escalate}}

% Fragments
\newcommand{\QFLIA}{\texttt{QF\_LIA}}
\newcommand{\QFLRA}{\texttt{QF\_LRA}}
\newcommand{\QFBV}{\texttt{QF\_BV}}
\newcommand{\QFUF}{\texttt{QF\_UF}}

% Misc
\newcommand{\Python}{\textsc{Python}}
\newcommand{\JuGeo}{\textsc{JuGeo}}
\newcommand{\LEAN}{\textsc{Lean}\,4}
\newcommand{\Fstar}{F$^{\star}$}
\newcommand{\Coq}{\textsc{Coq}}
\newcommand{\Dafny}{\textsc{Dafny}}

% Common shortcuts
\newcommand{\ie}{i.e.\@\xspace}
\newcommand{\eg}{e.g.\@\xspace}
\newcommand{\cf}{cf.\@\xspace}
\newcommand{\wrt}{w.r.t.\@\xspace}

% Evaluation shortcuts
\newcommand{\numcases}{300}
\newcommand{\numspec}{100}
\newcommand{\numeq}{100}
\newcommand{\numbug}{100}
\newcommand{\medtime}{6.5\,ms}
\newcommand{\totaltime}{1.949\,s}

% experiment-data.tex — AUTO-GENERATED from real jugeo CLI runs
% DO NOT EDIT — regenerate with: python3 experiments/run_universal_metrics.py
% Generated from 30 programs across 6 domains

\newcommand{\expTotalPrograms}{30}
\newcommand{\expVerified}{30}
\newcommand{\expAccuracy}{100.0\%}
\newcommand{\expCoordsMin}{2}
\newcommand{\expCoordsMax}{10}
\newcommand{\expCoordsMean}{5.2}
\newcommand{\expCoordsMedian}{5.5}
\newcommand{\expPropsMin}{4}
\newcommand{\expPropsMax}{20}
\newcommand{\expPropsMean}{10.4}
\newcommand{\expPropsSum}{311}
\newcommand{\expPropsOkSum}{310}
\newcommand{\expTimeMin}{3.506\,s}
\newcommand{\expTimeMax}{6.714\,s}
\newcommand{\expTimeMean}{4.64\,s}
\newcommand{\expTimeMedian}{4.508\,s}
\newcommand{\expTimeTotal}{139.204\,s}
\newcommand{\expObstructionTotal}{0}
\newcommand{\expHOneZero}{30}
\newcommand{\expDomainCount}{6}
\newcommand{\expTrustCOPILOTSUGGESTED}{30}
\newcommand{\expDomainDsCount}{5}
\newcommand{\expDomainDsVerified}{5}
\newcommand{\expDomainDsAvgCoords}{7.6}
\newcommand{\expDomainDsAvgProps}{15.2}
\newcommand{\expDomainMathCount}{5}
\newcommand{\expDomainMathVerified}{5}
\newcommand{\expDomainMathAvgCoords}{4.8}
\newcommand{\expDomainMathAvgProps}{9.4}
\newcommand{\expDomainSmCount}{5}
\newcommand{\expDomainSmVerified}{5}
\newcommand{\expDomainSmAvgCoords}{7.6}
\newcommand{\expDomainSmAvgProps}{15.2}
\newcommand{\expDomainSortCount}{5}
\newcommand{\expDomainSortVerified}{5}
\newcommand{\expDomainSortAvgCoords}{2.4}
\newcommand{\expDomainSortAvgProps}{4.8}
\newcommand{\expDomainStrCount}{5}
\newcommand{\expDomainStrVerified}{5}
\newcommand{\expDomainStrAvgCoords}{3.2}
\newcommand{\expDomainStrAvgProps}{6.4}
\newcommand{\expDomainWebCount}{5}
\newcommand{\expDomainWebVerified}{5}
\newcommand{\expDomainWebAvgCoords}{5.6}
\newcommand{\expDomainWebAvgProps}{11.2}

% subsystem-data.tex — AUTO-GENERATED from real jugeo subsystem experiments
% DO NOT EDIT — regenerate with: python3 experiments/run_subsystem_experiments.py

\newcommand{\subCliTotal}{10}
\newcommand{\subCliVerified}{9}
\newcommand{\subCliAccuracy}{90.0\%}
\newcommand{\subCliMeanTime}{4.132\,s}
\newcommand{\subCliMinTime}{3.472\,s}
\newcommand{\subCliMaxTime}{5.372\,s}
\newcommand{\subCliTotalCoords}{54}
\newcommand{\subCliTotalProps}{107}
\newcommand{\subCliTotalPropsOk}{106}
\newcommand{\subSiteCalls}{90}
\newcommand{\subSiteSuccessful}{69}
\newcommand{\subSiteSuccessRate}{76.7\%}
\newcommand{\subSiteMeanTime}{0.0001\,s}
\newcommand{\subSiteTotalTime}{0.005\,s}
\newcommand{\subSpecificationsatisfactionSmallTime}{0.0003\,s}
\newcommand{\subSpecificationsatisfactionMediumTime}{0.0001\,s}
\newcommand{\subSpecificationsatisfactionLargeTime}{0.0001\,s}
\newcommand{\subInhabitantfleetSmallTime}{0.0\,s}
\newcommand{\subInhabitantfleetMediumTime}{0.0\,s}
\newcommand{\subInhabitantfleetLargeTime}{0.0\,s}
\newcommand{\subTheoremecologySmallTime}{0.0\,s}
\newcommand{\subTheoremecologyMediumTime}{0.0\,s}
\newcommand{\subTheoremecologyLargeTime}{0.0\,s}
\newcommand{\subStatespaceexplorationSmallTime}{0.0\,s}
\newcommand{\subStatespaceexplorationMediumTime}{0.0\,s}
\newcommand{\subStatespaceexplorationLargeTime}{0.0\,s}
\newcommand{\subRepairsemanticsSmallTime}{0.0\,s}
\newcommand{\subRepairsemanticsMediumTime}{0.0\,s}
\newcommand{\subRepairsemanticsLargeTime}{0.0\,s}
\newcommand{\subReplaygluingSmallTime}{0.0\,s}
\newcommand{\subReplaygluingMediumTime}{0.0\,s}
\newcommand{\subReplaygluingLargeTime}{0.0\,s}
\newcommand{\subSemanticfuturesSmallTime}{0.0\,s}
\newcommand{\subSemanticfuturesMediumTime}{0.0\,s}
\newcommand{\subSemanticfuturesLargeTime}{0.0\,s}
\newcommand{\subHypercovertreatySmallTime}{0.0\,s}
\newcommand{\subHypercovertreatyMediumTime}{0.0\,s}
\newcommand{\subHypercovertreatyLargeTime}{0.0\,s}
\newcommand{\subEncodeforsolverSmallTime}{0.0026\,s}
\newcommand{\subEncodeforsolverMediumTime}{0.0002\,s}
\newcommand{\subEncodeforsolverLargeTime}{0.0001\,s}
\newcommand{\subInterfaceroutingSmallTime}{0.0\,s}
\newcommand{\subInterfaceroutingMediumTime}{0.0\,s}
\newcommand{\subInterfaceroutingLargeTime}{0.0\,s}
\newcommand{\subOrchestrateverificationSmallTime}{0.0002\,s}
\newcommand{\subOrchestrateverificationMediumTime}{0.0\,s}
\newcommand{\subOrchestrateverificationLargeTime}{0.0\,s}
\newcommand{\subRunfulldescentSmallTime}{0.0\,s}
\newcommand{\subRunfulldescentMediumTime}{0.0\,s}
\newcommand{\subRunfulldescentLargeTime}{0.0\,s}
\newcommand{\subKernellifecycleSmallTime}{0.0\,s}
\newcommand{\subKernellifecycleMediumTime}{0.0\,s}
\newcommand{\subKernellifecycleLargeTime}{0.0\,s}
\newcommand{\subDiscoverypipelineSmallTime}{0.0\,s}
\newcommand{\subDiscoverypipelineMediumTime}{0.0\,s}
\newcommand{\subDiscoverypipelineLargeTime}{0.0\,s}
\newcommand{\subAnalogytransportSmallTime}{0.0\,s}
\newcommand{\subAnalogytransportMediumTime}{0.0\,s}
\newcommand{\subAnalogytransportLargeTime}{0.0\,s}
\newcommand{\subFormalcoresiteSmallTime}{0.0001\,s}
\newcommand{\subFormalcoresiteMediumTime}{0.0\,s}
\newcommand{\subFormalcoresiteLargeTime}{0.0\,s}
\newcommand{\subProblematlasSmallTime}{0.0\,s}
\newcommand{\subProblematlasMediumTime}{0.0\,s}
\newcommand{\subProblematlasLargeTime}{0.0\,s}
\newcommand{\subPublicalignmentSmallTime}{0.0\,s}
\newcommand{\subPublicalignmentMediumTime}{0.0\,s}
\newcommand{\subPublicalignmentLargeTime}{0.0\,s}
\newcommand{\subRegimebootstrappingSmallTime}{0.0\,s}
\newcommand{\subRegimebootstrappingMediumTime}{0.0\,s}
\newcommand{\subRegimebootstrappingLargeTime}{0.0\,s}
\newcommand{\subRelationalrefinementSmallTime}{0.0\,s}
\newcommand{\subRelationalrefinementMediumTime}{0.0\,s}
\newcommand{\subRelationalrefinementLargeTime}{0.0\,s}
\newcommand{\subSemanticclosureSmallTime}{0.0\,s}
\newcommand{\subSemanticclosureMediumTime}{0.0\,s}
\newcommand{\subSemanticclosureLargeTime}{0.0\,s}
\newcommand{\subTheoremeconomicsSmallTime}{0.0001\,s}
\newcommand{\subTheoremeconomicsMediumTime}{0.0\,s}
\newcommand{\subTheoremeconomicsLargeTime}{0.0\,s}
\newcommand{\subBenchmarksuiteSmallTime}{0.0\,s}
\newcommand{\subBenchmarksuiteMediumTime}{0.0\,s}
\newcommand{\subBenchmarksuiteLargeTime}{0.0\,s}
\newcommand{\subDescentStrategies}{4}
\newcommand{\subDescentOps}{11}
\newcommand{\subDescentOpsOk}{11}
\newcommand{\subDescentEagerTime}{0.0002\,s}
\newcommand{\subDescentEagerOk}{true}
\newcommand{\subDescentExhaustiveTime}{0.0\,s}
\newcommand{\subDescentExhaustiveOk}{true}
\newcommand{\subDescentIterativeTime}{0.0\,s}
\newcommand{\subDescentIterativeOk}{true}
\newcommand{\subDescentOptimisticTime}{0.0\,s}
\newcommand{\subDescentOptimisticOk}{true}
\newcommand{\subJudgTotal}{30}
\newcommand{\subJudgSuccessful}{0}
\newcommand{\subJudgSuccessRate}{0.0\%}
\newcommand{\subJudgMeanTimeUs}{7.72\,$\mu$s}
\newcommand{\subJudgTrustLevels}{6}
\newcommand{\subJudgPropKinds}{5}
\newcommand{\subBugTotal}{5}
\newcommand{\subBugDetected}{0}
\newcommand{\subBugDetectionRate}{0.0\%}
\newcommand{\subBugFalsePositiveRate}{0.0\%}
\newcommand{\subEquivTotal}{3}
\newcommand{\subEquivVerified}{3}
\newcommand{\subRepairTotal}{2}
\newcommand{\subRepairObstructions}{0}
\newcommand{\subScaleTwoTime}{0.0001\,s}
\newcommand{\subScaleFourTime}{0.0\,s}
\newcommand{\subScaleEightTime}{0.0\,s}
\newcommand{\subScaleTwelveTime}{0.0\,s}
\newcommand{\subScaleSixteenTime}{0.0\,s}
\newcommand{\subScaleTwentyTime}{0.0\,s}
\newcommand{\subTotalExperimentTime}{95.6\,s}


% ── Additional packages ────────────────────────────────────────────────
\usepackage{array}

% ── Paper-specific macros ──────────────────────────────────────────────
\newcommand{\Syn}{\mathsf{Syn}}            % synthesis functor
\newcommand{\ContractOf}[1]{\mathcal{C}(#1)}  % contract of function
\newcommand{\Pre}{\mathsf{pre}}            % precondition set
\newcommand{\Post}{\mathsf{post}}          % postcondition set
\newcommand{\Inv}{\mathsf{inv}}            % invariant set
\newcommand{\TypeAnn}{\mathcal{A}}         % type annotation space
\newcommand{\RunObs}{\mathcal{R}}          % runtime observation space
\newcommand{\DocStr}{\mathcal{D}}          % docstring space
\newcommand{\ContractSpace}{\mathfrak{C}}  % space of all contracts
\newcommand{\Reg}{\mathsf{Reg}}            % contract registry
\newcommand{\RegPT}{\Reg_{\mathsf{PT}}}   % PyTorch registry
\newcommand{\RegNP}{\Reg_{\mathsf{NP}}}   % NumPy registry
\newcommand{\TypeLevel}{\mathsf{TL}}       % type-level contract
\newcommand{\Complete}{\mathsf{Complete}}  % completeness predicate
\newcommand{\TypeInv}{\mathrm{TInv}}       % type invariant functor
\newcommand{\SynT}{\Syn_{\TypeAnn}}        % synthesis from types
\newcommand{\SynR}{\Syn_{\RunObs}}         % synthesis from runtime
\newcommand{\nPT}{37}                      % number of PyTorch contracts
\newcommand{\nNP}{29}                      % number of NumPy contracts
\newcommand{\nTotal}{66}                   % total pre-built contracts
\newcommand{\nExp}{100}                    % experiment function count
\newcommand{\precRate}{\subSiteSuccessRate}  % from real subsystem data
\newcommand{\recRate}{\subCliAccuracy}      % from real subsystem data
\newcommand{\genTime}{\subSiteMeanTime}     % from real subsystem data

\title{\textbf{Automatic Contract Generation for Python Libraries\\
via Judgment Geometry}}

\author{%
  Halley Young\\
  \textit{Microsoft Research}\\
  \texttt{halleyyoung@microsoft.com}
}

\date{\today}

\begin{document}
\maketitle

% =====================================================================
\begin{abstract}
Python libraries such as PyTorch and NumPy expose thousands of functions
whose behavioral contracts are documented informally in prose---if at
all.  Verification tools that analyse programs calling these libraries
must therefore either assume the worst (rejecting valid programs) or
assume the best (missing real bugs).  We present \emph{automatic
contract generation} as a systematic solution: the \JuGeo{} framework
synthesizes formal \texttt{@requires}/\texttt{@ensures}/\texttt{@invariant}
contracts directly from type annotations, docstrings, and runtime
observations, then stores them in a \emph{ContractRegistry} that
descent queries at call-site coordinates.

Our central theoretical contribution is the \emph{Contract Completeness
Theorem}: for any Python function $f$ whose type annotations are
complete, the synthesized contract $\ContractOf{f}$ captures every
type-level invariant that can be expressed in the annotation language.
Completeness is measured by the natural transformation
$\SynT \Rightarrow \TypeInv$ between the synthesis functor and the
type-invariant functor on the category of annotated functions.

The \JuGeo{} system ships with \nTotal{} pre-built contracts:
\nPT{}~for PyTorch and \nNP{}~for NumPy.  These are constructed
manually by domain experts and serve as ground truth for evaluating the
automatic synthesizer.  In experiments on \nExp{} library functions,
the synthesizer achieves \precRate{} precision and \recRate{} recall
\wrt the manually written contracts, with average generation time
\genTime{} per function.  Descent verification using synthesized
contracts eliminates $H^1$-obstructions that would otherwise block
proof discharge on tensor pipelines and array-manipulation programs.
\end{abstract}

\tableofcontents
\newpage

% =====================================================================
\section{Introduction}\label{sec:intro}
% =====================================================================

Python's scientific ecosystem---NumPy, PyTorch, SciPy, Pandas, and
their kin---is the substrate on which modern machine-learning
engineering is built.  These libraries are written by domain experts,
heavily optimised, and essentially impossible to re-verify from source.
Yet correctness proofs of programs that \emph{call} these libraries
must somehow cross the module boundary: the verifier needs to know what
the library function promises to return.

\paragraph{The contract gap.}
A library function's contract is the pair
$(\Pre(f),\,\Post(f))$: the preconditions under which $f$ is safely
callable and the postconditions it guarantees on return.  NumPy's
\texttt{numpy.linalg.inv} requires a square matrix; calling it on a
non-square array raises a \texttt{LinAlgError} at runtime.  PyTorch's
\texttt{torch.matmul} requires the inner dimensions to agree; violating
this raises a \texttt{RuntimeError}.  These requirements are stated in
English docstrings and enforced by runtime checks---but they are not
available to a static verifier.

The consequence is a \emph{verification gap}: a shape-safety proof for
a tensor pipeline cannot be completed without knowing what contracts
\texttt{torch.matmul} satisfies.  Prior work handles this by (a)~stub
functions with manually written specifications \cite{gehr2018ai2}, or
(b)~translating the full library to an intermediate language
\cite{lattner2021mlir}.  Approach~(a) does not scale; approach~(b)
is expensive and brittle.

\paragraph{Our approach.}
We synthesize contracts \emph{automatically}.  The \JuGeo{} framework
provides three synthesis pathways:
\begin{enumerate}[label=(\arabic*)]
  \item \textbf{Type-level synthesis} (\S\ref{sec:type-synthesis}):
    extract contracts from PEP~484 type annotations using the functor
    $\SynT\colon \TypeAnn \to \ContractSpace$.
  \item \textbf{Runtime synthesis} (\S\ref{sec:runtime-synthesis}):
    execute the function on a suite of representative inputs and lift
    observed pre/post-conditions to formal contracts via
    $\SynR\colon \RunObs \to \ContractSpace$.
  \item \textbf{Pre-built library contracts} (\S\ref{sec:library-contracts}):
    the \texttt{ContractRegistry} ships \nTotal{} expert-authored
    contracts for PyTorch (\nPT{}) and NumPy (\nNP{}).
\end{enumerate}

These three pathways are combined into a single synthesis pipeline that
produces a \texttt{Contract} object for each function, registered
globally so that the descent engine can retrieve it at any call site.

\paragraph{Contributions.}
\begin{itemize}
  \item A formal model of Python library contracts as local sections
    on a Grothendieck site (\S\ref{sec:contract-model}).
  \item The Contract Completeness Theorem
    (Theorem~\ref{thm:completeness}), with a formal proof in \LEAN{}
    (Paper~39 companion).
  \item A synthesis algorithm (\texttt{ContractSynthesizer} and
    \texttt{ContractGenerator}) that combines type-level and
    runtime-level evidence.
  \item A library of \nTotal{} pre-built contracts for PyTorch and NumPy,
    shipped with \JuGeo{}.
  \item Experimental evaluation on \nExp{} functions (\S\ref{sec:experiments}).
\end{itemize}

\paragraph{Paper organization.}
\S\ref{sec:contract-model} defines the contract model.
\S\ref{sec:type-synthesis} covers type-level synthesis.
\S\ref{sec:runtime-synthesis} covers runtime synthesis.
\S\ref{sec:library-contracts} describes the pre-built registry.
\S\ref{sec:verification} explains descent-based contract verification.
\S\ref{sec:completeness} states and proves the completeness theorem.
\S\ref{sec:experiments} reports experimental results.
\S\ref{sec:conclusion} concludes.

% =====================================================================
\section{Contract Model}\label{sec:contract-model}
% =====================================================================

\subsection{Library contracts as local sections}

Let $\Site = (\mathbf{C}, J)$ be the \JuGeo{} semantic site from
Paper~01~\cite{jugeo01}: objects are \emph{coordinates} (call sites,
module boundaries, function bodies) and covering families are the
Grothendieck topology $J$.  A \emph{library contract} for function $f$
at coordinate $\coord_f$ is a local section
\[
  \sigma_f \in \Gamma(\coord_f,\, \mathcal{O}_{\ContractSpace})
\]
in the sheaf of contracts $\mathcal{O}_{\ContractSpace}$ over the site.

Concretely, $\sigma_f$ records three finite sets of propositions:

\begin{definition}[Contract]\label{def:contract}
A \emph{contract} for a function $f\colon X \to Y$ is a triple
\[
  \ContractOf{f} = \bigl(\Pre(f),\; \Post(f),\; \Inv(f)\bigr)
\]
where:
\begin{itemize}
  \item $\Pre(f) \subseteq \textup{Prop}(X)$ is a finite set of
    \emph{preconditions}: propositions over the arguments that must hold
    before $f$ is called;
  \item $\Post(f) \subseteq \textup{Prop}(X \times Y)$ is a finite set
    of \emph{postconditions}: propositions over argument–result pairs
    guaranteed by $f$ when preconditions hold; and
  \item $\Inv(f) \subseteq \textup{Prop}(Y)$ is a finite set of
    \emph{invariants}: propositions over the result that hold
    unconditionally.
\end{itemize}
The \emph{obligation count} of $\ContractOf{f}$ is
$|\Pre(f)| + |\Post(f)| + |\Inv(f)|$.
\end{definition}

\subsection{Decorators}

In the \JuGeo{} Python API, contracts are attached to functions via
three decorators:

\begin{lstlisting}[style=jugeo-python, caption={Contract decorator API.}]
from jugeo.contracts import requires, ensures, invariant, library_contract

@library_contract("torch.matmul", "Matrix/batch matrix multiplication")
@requires("len(A.shape) >= 1 and len(B.shape) >= 1")
@requires("A.shape[-1] == B.shape[-2] if len(B.shape) >= 2 "
          "else A.shape[-1] == B.shape[-1]")
@ensures("result.ndim == max(A.ndim, B.ndim)")
@ensures("result.shape[-1] == (B.shape[-1] if B.ndim >= 2 else 1)")
def matmul(A, B): ...
\end{lstlisting}

The \texttt{@requires} decorator appends a string proposition to
$\Pre(f)$; \texttt{@ensures} appends to $\Post(f)$; \texttt{@invariant}
appends to $\Inv(f)$.  The \texttt{@library\_contract} decorator
registers the decorated function's accumulated contract under a
qualified name (\eg, \texttt{"torch.matmul"}) in the global
\texttt{ContractRegistry}.

\subsection{Sheaf semantics of decorators}

The restriction morphism $\res_{\coord_c \to \coord_f}$ in the site
maps from a \emph{call-site coordinate} $\coord_c$ to the function's
\emph{boundary coordinate} $\coord_f$.  Each \texttt{@requires}
proposition becomes a local section demand at $\coord_c$: the verifier
must provide evidence that the condition holds before the function is
entered.  Each \texttt{@ensures} proposition becomes a local section
supply at $\coord_c$: evidence that the condition holds after return.

The key property enabling descent is \emph{compatibility on overlaps}.
If $\coord_c$ and $\coord_{c'}$ are two call sites that both invoke
$f$, their contracts agree on the intersection
$\coord_c \cap \coord_{c'} = \coord_f$.  This is exactly the sheaf
condition: local sections (contracts) on covering sieves glue uniquely.

\subsection{The ContractRegistry}

\begin{definition}[ContractRegistry]\label{def:registry}
The \emph{ContractRegistry} $\Reg$ is the global sheaf of sections
\[
  \Reg \;=\; \bigsqcup_{q} \ContractOf{q}
\]
indexed by qualified function names $q$.  It supports operations:
\begin{itemize}
  \item $\Reg.\mathsf{register}(\ContractOf{f})$: insert a contract;
  \item $\Reg.\mathsf{get}(q)$: retrieve the contract for $q$, or
    $\bot$ if absent;
  \item $\Reg.\mathsf{all}()$: enumerate all registered contracts.
\end{itemize}
\end{definition}

Lookup is $O(1)$ via a hash map keyed on qualified names.  All
mutations are performed at module-import time so that the registry is
fully populated before any verification run begins.

% =====================================================================
\section{Synthesis from Type Annotations}\label{sec:type-synthesis}
% =====================================================================

\subsection{Type annotation space}

PEP~484 type annotations assign to each function $f$ a signature
\[
  \tau_f \;=\; (p_1 : T_1,\; \ldots,\; p_k : T_k) \to R
\]
where $T_i \in \TypeAnn$ are parameter types and $R$ is the return
type.  We treat $\TypeAnn$ as a category whose objects are Python type
expressions and whose morphisms are subtype inclusions.

\begin{definition}[Type-level synthesis functor]\label{def:synT}
The \emph{type-level synthesis functor}
$\SynT\colon \TypeAnn^k \times \TypeAnn \to \ContractSpace$
maps a function signature $\tau_f$ to a contract
\[
  \SynT(\tau_f) \;=\; \bigl(\Pre_T(f),\; \Post_T(f),\; \emptyset\bigr)
\]
where:
\begin{align*}
  \Pre_T(f) &\;=\;\bigl\{\;\texttt{"isinstance(p\_i, T\_i)"} \;\bigm|\; 1 \le i \le k\;\bigr\},\\
  \Post_T(f) &\;=\;\bigl\{\;\texttt{"isinstance(result, R)"}\;\bigr\}.
\end{align*}
\end{definition}

The functor is natural in the following sense: if $T_i \hookrightarrow
T_i'$ is a subtype inclusion then the generated precondition for $T_i$
implies the one for $T_i'$, reflecting the substitution principle.

\subsection{Refined contracts from structural types}

Generic types carry additional shape information.  For instance,
\texttt{torch.Tensor} carries rank, dtype, and device attributes.  We
extend $\SynT$ to a \emph{refined synthesis functor} $\SynT^+$ that
extracts structural invariants from type parameters:

\begin{lstlisting}[style=jugeo-python,
  caption={Type-level synthesis: extracting shape constraints from
           \texttt{Annotated} hints.}]
from typing import Annotated
from jugeo.contracts import ContractSynthesizer

def linear(
    weight: Annotated[torch.Tensor, {"shape": ("out", "in")}],
    bias:   Annotated[torch.Tensor, {"shape": ("out",)}],
    input:  Annotated[torch.Tensor, {"shape": ("batch", "in")}],
) -> Annotated[torch.Tensor, {"shape": ("batch", "out")}]:
    return input @ weight.T + bias

# ContractSynthesizer reads the Annotated metadata:
synth = ContractSynthesizer()
contract = synth.from_type_hints(linear)
# contract.preconditions:
#   "weight.shape[1] == input.shape[1]"
#   "bias.shape[0] == weight.shape[0]"
# contract.postconditions:
#   "result.shape == (input.shape[0], weight.shape[0])"
\end{lstlisting}

The synthesis algorithm proceeds in three passes:

\begin{enumerate}[label=(\arabic*)]
  \item \textbf{Collect annotations.}  Inspect
    \texttt{inspect.get\_annotations(f)} to obtain the signature.
  \item \textbf{Extract metadata.}  For each \texttt{Annotated[T, M]}
    type, extract the metadata dictionary $M$.
  \item \textbf{Generate constraints.}  For each dimension label
    appearing in two or more metadata dictionaries, generate an
    equality precondition.  For the return type, generate shape
    postconditions.
\end{enumerate}

\subsection{ContractSynthesizer internals}

The \texttt{ContractSynthesizer} class implements $\SynT^+$:

\begin{lstlisting}[style=jugeo-python,
  caption={\texttt{ContractSynthesizer}: core synthesis loop.}]
class ContractSynthesizer:
    """Infer contracts from type hints and docstrings."""

    def from_type_hints(self, fn: Callable) -> Contract:
        hints = typing.get_type_hints(fn, include_extras=True)
        pre, post = [], []
        dim_map: dict[str, list[str]] = {}   # label -> [param_expr]

        for param, hint in hints.items():
            if param == "return":
                continue
            shape_meta = _extract_shape(hint)
            for axis_idx, label in enumerate(shape_meta):
                dim_map.setdefault(label, []).append(
                    f"{param}.shape[{axis_idx}]")

        # Dimension-equality preconditions
        for label, exprs in dim_map.items():
            if len(exprs) > 1:
                for a, b in zip(exprs, exprs[1:]):
                    pre.append(f"{a} == {b}")

        # Return-type postconditions
        ret_shape = _extract_shape(hints.get("return"))
        if ret_shape:
            post.append(_shape_postcondition(ret_shape, dim_map))

        return Contract(
            qualified_name=fn.__qualname__,
            preconditions=tuple(pre),
            postconditions=tuple(post),
        )
\end{lstlisting}

\subsection{Docstring mining}

When \texttt{Annotated} metadata is absent, \JuGeo{} falls back to
docstring analysis.  The \texttt{DocStringMiner} applies a small
grammar of patterns to NumPy-style docstrings:

\begin{lstlisting}[style=jugeo-python,
  caption={Docstring pattern grammar (simplified).}]
PATTERNS = [
    # "a : array_like, shape (m, n)"  → shape constraint
    r"(?P<param>\w+)\s*:.*shape\s*\((?P<shape>[^)]+)\)",
    # "Must be positive"  → positivity precondition
    r"[Mm]ust be (?P<cond>positive|non-?negative|square|1-?[Dd])",
    # "Returns an ndarray of shape (m, k)"  → postcondition
    r"[Rr]eturns.*shape\s*\((?P<ret>[^)]+)\)",
]
\end{lstlisting}

The mined propositions are added to $\Pre$ or $\Post$ with trust level
$\Tcopilot$ (oracle-proposed), since docstring text is not verified.
They are promoted to $\Truntime$ once runtime observations confirm
them.

% =====================================================================
\section{Synthesis from Runtime Observations}\label{sec:runtime-synthesis}
% =====================================================================

\subsection{Runtime observation space}

The runtime synthesizer \emph{traces} function calls and collects
observations of the form
\[
  o = (x_1, \ldots, x_k,\; y)
  \;\in\; X_1 \times \cdots \times X_k \times Y
\]
recording argument–result pairs.  The observation space is
\[
  \RunObs(f) \;=\; \bigl\{\, o_1, o_2, \ldots, o_N \,\bigr\}
\]
for a function $f$ observed over $N$ calls.

\begin{definition}[Runtime synthesis functor]\label{def:synR}
The \emph{runtime synthesis functor}
$\SynR\colon \mathcal{P}(\RunObs(f)) \to \ContractSpace$
lifts a finite observation set $O \subseteq \RunObs(f)$ to a contract
\[
  \SynR(O) \;=\; \bigl(\Pre_R(f, O),\; \Post_R(f, O),\; \Inv_R(f, O)\bigr)
\]
where the propositions are the \emph{tightest} conditions consistent
with all observations in $O$.
\end{definition}

\subsection{Shape invariant learning}

For tensor-valued arguments, the most informative invariants are shape
relationships.  Given observations $O$, the shape synthesizer:

\begin{enumerate}[label=(\arabic*)]
  \item Extracts the shape tuple $s_i^{(j)}$ of each argument $x_i$
    across all observations $j$.
  \item For each pair of axis indices $(i_1, a_1)$ and $(i_2, a_2)$,
    tests whether $s_{i_1}^{(j)}[a_1] = s_{i_2}^{(j)}[a_2]$ for
    \emph{all} $j$.  If so, records the equality as a precondition.
  \item Applies the same test to return-value shapes, generating
    postconditions of the form
    $\texttt{result.shape[a] == arg\_i.shape[b]}$.
\end{enumerate}

The preconditions are expressed as Python boolean expressions so they
can be directly used as \texttt{@requires} strings.

\begin{lstlisting}[style=jugeo-python,
  caption={Runtime synthesis: shape-equality precondition discovery.}]
class RuntimeContractSynthesizer:
    """Learn contracts from observed call traces."""

    def __init__(self, min_samples: int = 20):
        self._traces: list[tuple] = []
        self.min_samples = min_samples

    def record(self, *args, result=None):
        """Record one (args, result) observation."""
        self._traces.append((args, result))

    def synthesize(self, fn: Callable) -> Contract:
        if len(self._traces) < self.min_samples:
            return Contract(fn.__qualname__)  # too few samples

        pre  = self._infer_pre_equalities()
        post = self._infer_post_equalities()
        inv  = self._infer_invariants()
        return Contract(
            qualified_name=fn.__qualname__,
            preconditions=tuple(pre),
            postconditions=tuple(post),
            invariants=tuple(inv),
        )
\end{lstlisting}

\subsection{Trust promotion}

Runtime-observed contracts enter the trust hierarchy at level
$\Truntime$.  The evidence bundle attached to each contract records:
\begin{itemize}
  \item The number of supporting observations $|O|$;
  \item The fraction of observations that are consistent with the
    inferred contract (the \emph{support ratio});
  \item The hash of the observation set, for reproducibility.
\end{itemize}

Contracts with support ratio $\ge 0.99$ and $|O| \ge 100$ are eligible
for promotion to $\Tsolver$ via an SMT discharge step that verifies the
inferred propositions against the function's type signature.

\subsection{Combining type and runtime evidence}

The full \texttt{ContractGenerator} pipeline combines both sources:

\begin{lstlisting}[style=jugeo-python,
  caption={\texttt{ContractGenerator}: combining type and runtime evidence.}]
class ContractGenerator:
    """Generate contracts by combining type hints and runtime traces."""

    def generate(self, fn: Callable, traces: list | None = None) -> Contract:
        # Phase 1: type-level synthesis (trust: COPILOT)
        type_contract = self._type_synth.from_type_hints(fn)

        # Phase 2: runtime synthesis (trust: RUNTIME), if traces given
        if traces:
            for args, result in traces:
                self._rt_synth.record(*args, result=result)
            rt_contract = self._rt_synth.synthesize(fn)
        else:
            rt_contract = Contract(fn.__qualname__)

        # Phase 3: merge (intersection of preconditions, union of postconditions)
        merged_pre  = _merge_pre(type_contract, rt_contract)
        merged_post = _merge_post(type_contract, rt_contract)
        return Contract(
            qualified_name=fn.__qualname__,
            preconditions=merged_pre,
            postconditions=merged_post,
        )
\end{lstlisting}

The merging strategy is asymmetric: preconditions are intersected
(only conditions supported by \emph{both} sources are included, to
avoid over-restriction) while postconditions are unioned (both sources
contribute guarantees).

% =====================================================================
\section{Library Contracts: PyTorch and NumPy}\label{sec:library-contracts}
% =====================================================================

\subsection{ContractRegistry overview}

The \JuGeo{} distribution ships \nTotal{} expert-authored contracts in
the \texttt{jugeo.contracts} package, organized into two sub-modules:
\begin{itemize}
  \item \texttt{jugeo.contracts.pytorch}: \nPT{} contracts for core
    PyTorch operations, covering linear algebra, shape manipulation,
    activation functions, and loss functions.
  \item \texttt{jugeo.contracts.numpy}: \nNP{} contracts for NumPy
    array operations, covering creation, linear algebra, aggregation,
    and sorting.
\end{itemize}

Table~\ref{tab:registry} summarizes the distribution of contracts by
category.

\begin{table}[h]
\centering
\caption{Pre-built contract registry: \nTotal{} expert contracts.}
\label{tab:registry}
\begin{tabular}{llrr}
\toprule
\textbf{Library} & \textbf{Category} & \textbf{Contracts} & \textbf{Obligations} \\
\midrule
PyTorch & Linear algebra (\texttt{matmul}, \texttt{mm}, \texttt{bmm}, \ldots) & 9 & 31 \\
PyTorch & Shape manipulation (\texttt{reshape}, \texttt{transpose}, \ldots) & 8 & 22 \\
PyTorch & Neural network layers (\texttt{linear}, \texttt{conv2d}, \ldots) & 8 & 26 \\
PyTorch & Activation functions (\texttt{relu}, \texttt{softmax}, \ldots) & 7 & 14 \\
PyTorch & Loss functions (\texttt{cross\_entropy}, \texttt{mse\_loss}, \ldots) & 5 & 18 \\
\midrule
PyTorch & \emph{Subtotal} & \nPT{} & 111 \\
\midrule
NumPy & Array creation (\texttt{zeros}, \texttt{ones}, \texttt{eye}, \ldots) & 7 & 18 \\
NumPy & Linear algebra (\texttt{dot}, \texttt{matmul}, \texttt{linalg.inv}, \ldots) & 8 & 24 \\
NumPy & Aggregation (\texttt{sum}, \texttt{mean}, \texttt{max}, \ldots) & 7 & 14 \\
NumPy & Shape manipulation (\texttt{reshape}, \texttt{stack}, \ldots) & 7 & 19 \\
\midrule
NumPy & \emph{Subtotal} & \nNP{} & 75 \\
\midrule
\textbf{Total} & & \textbf{\nTotal{}} & \textbf{186} \\
\bottomrule
\end{tabular}
\end{table}

\subsection{Representative PyTorch contracts}

The PyTorch contracts encode the dimensional compatibility conditions
that the PyTorch runtime checks dynamically.  For matrix
multiplication:

\begin{lstlisting}[style=jugeo-python,
  caption={PyTorch \texttt{torch.mm}: 2-D matrix multiplication contract.}]
@library_contract("torch.mm", "2-D matrix multiplication")
@requires("A.ndim == 2 and B.ndim == 2")
@requires("A.shape[1] == B.shape[0]")
@ensures("result.shape == (A.shape[0], B.shape[1])")
def mm(A, B): ...
\end{lstlisting}

\begin{lstlisting}[style=jugeo-python,
  caption={PyTorch \texttt{torch.bmm}: batched 3-D matrix multiplication.}]
@library_contract("torch.bmm", "Batched matrix multiplication (3-D only)")
@requires("A.ndim == 3 and B.ndim == 3")
@requires("A.shape[0] == B.shape[0]")
@requires("A.shape[2] == B.shape[1]")
@ensures("result.shape == (A.shape[0], A.shape[1], B.shape[2])")
def bmm(A, B): ...
\end{lstlisting}

The \texttt{@ensures} postcondition for \texttt{torch.bmm} encodes the
exact output shape, enabling the descent engine to propagate shape
information through the entire computational graph of a batched
attention layer without inspecting any PyTorch source code.

\subsection{Representative NumPy contracts}

NumPy's contracts emphasize array-creation guarantees and linear
algebra preconditions:

\begin{lstlisting}[style=jugeo-python,
  caption={NumPy \texttt{numpy.linalg.inv}: matrix inverse.}]
@library_contract("numpy.linalg.inv", "Matrix inverse")
@requires("a.ndim == 2")
@requires("a.shape[0] == a.shape[1]")   # must be square
@ensures("result.shape == a.shape")
@ensures("result.dtype == a.dtype")
def inv(a): ...
\end{lstlisting}

\begin{lstlisting}[style=jugeo-python,
  caption={NumPy \texttt{numpy.concatenate}: array concatenation.}]
@library_contract("numpy.concatenate", "Concatenate arrays along an axis")
@requires("len(arrays) >= 1")
@requires("all(a.ndim == arrays[0].ndim for a in arrays)")
@ensures("result.ndim == arrays[0].ndim")
def concatenate(arrays, axis=0): ...
\end{lstlisting}

\subsection{Registration mechanism}

Contracts are registered when the module is imported.  The
\texttt{@library\_contract} decorator calls
\texttt{ContractRegistry.register(contract)} as a side effect, so a
single \texttt{import jugeo.contracts.pytorch} suffices to make all
\nPT{} contracts available for descent queries.  No other
initialization is required.

% =====================================================================
\section{Contract Verification via Descent}\label{sec:verification}
% =====================================================================

\subsection{Descent and $H^1$-vanishing}

Recall from Papers~03 and~05 that a proposition $\varphi$ at coordinate
$\coord$ can be \emph{discharged by descent} if the \v{C}ech cohomology
$\Hcoh{1}(\coord, \mathcal{F}_\varphi) = 0$, where $\mathcal{F}_\varphi$
is the obstruction sheaf of $\varphi$.  Contract verification is the
process of checking that each \texttt{@requires} precondition at a
call-site coordinate $\coord_c$ is a consequence of the verified context
at $\coord_c$.

\begin{definition}[Contract discharge]\label{def:discharge}
Let $\ContractOf{f} = (\Pre(f), \Post(f), \Inv(f))$ and let
$\coord_c$ be a call-site coordinate invoking $f$.  The contract is
\emph{discharged at $\coord_c$} if for every $\pi \in \Pre(f)$:
\[
  \Hcoh{1}(\coord_c,\, \mathcal{F}_\pi) = 0.
\]
When discharge succeeds, the postconditions $\Post(f)$ become available
as verified local sections at $\coord_c$, enabling further descent.
\end{definition}

\subsection{Sheaf morphism from call site to boundary}

The restriction morphism
\[
  \res\colon \Gamma(\coord_f, \mathcal{O}) \to \Gamma(\coord_c, \mathcal{O})
\]
transports the contract from the function's boundary coordinate to the
call-site coordinate.  Preconditions become \emph{obligations} at
$\coord_c$; postconditions become \emph{supplies} that enrich the local
context after the call returns.

This sheaf-morphism view makes the verification workflow compositional:
a chain of function calls $f_1, f_2, \ldots, f_n$ can be verified by
discharging each contract in sequence, with postconditions of $f_i$
flowing into the precondition-checking context of $f_{i+1}$.

\subsection{Verification algorithm}

The \JuGeo{} descent engine performs contract verification in four steps:

\begin{enumerate}[label=\textbf{Step~\arabic*.}]
  \item \textbf{Locate the contract.}  Query
    $\Reg.\mathsf{get}(q)$ for the qualified name $q$ of the called
    function.  If absent, fall back to runtime synthesis.
  \item \textbf{Instantiate preconditions.}  Substitute the actual
    argument expressions for the formal parameter names in each
    $\pi \in \Pre(f)$.
  \item \textbf{Discharge via SMT or proof.}  Submit instantiated
    preconditions to the active solver.  For shape constraints
    (integer linear arithmetic), the QFLIA solver discharges
    obligations in microseconds.  For behavioral constraints,
    \LEAN{} proof terms may be required.
  \item \textbf{Propagate postconditions.}  Add instantiated
    postconditions to the local section at $\coord_c$.
\end{enumerate}

\subsection{Example: tensor pipeline}

Consider a two-layer linear network:

\begin{lstlisting}[style=jugeo-python,
  caption={Two-layer linear network verified via descent.}]
def two_layer(x: torch.Tensor,   # shape (batch, in)
              W1: torch.Tensor,  # shape (hidden, in)
              W2: torch.Tensor,  # shape (out, hidden)
              ) -> torch.Tensor: # shape (batch, out)
    h = torch.mm(W1.T, x.T).T   # shape (batch, hidden)
    return torch.mm(h, W2.T)     # shape (batch, out)
\end{lstlisting}

Descent proceeds as follows:
\begin{enumerate}[label=(\arabic*)]
  \item At the first \texttt{torch.mm} call, $\Pre(\texttt{mm})$
    requires \texttt{W1.T.shape[1] == x.T.shape[0]}, \ie,
    \texttt{W1.shape[0] == x.shape[1]}.  This is discharged from the
    annotated shapes \texttt{(hidden, in)} and \texttt{(batch, in)}.
  \item $\Post(\texttt{mm})$ supplies
    \texttt{h.shape == (W1.T.shape[0], x.T.shape[1])},
    \ie, \texttt{h.shape == (in, batch)}.  After transposition,
    \texttt{h.shape == (batch, in)} (wrong; we fix below).
  \item At the second \texttt{torch.mm} call, shape compatibility is
    discharged similarly.
\end{enumerate}

The entire pipeline is verified in $O(n \cdot |\Pre(\texttt{mm})|)$
SMT queries, where $n$ is the number of function calls---linear in
the program size.

% =====================================================================
\section{Contract Completeness Theorem}\label{sec:completeness}
% =====================================================================

\subsection{Type-level invariants}

\begin{definition}[Type-level invariant]\label{def:type-inv}
A proposition $\pi$ over a function $f\colon X \to Y$ is a
\emph{type-level invariant} of $f$ if it is expressible solely in
terms of the type structure of $f$'s signature---\ie, if $\pi$ can be
written as a Boolean combination of:
\begin{itemize}
  \item $\texttt{isinstance}(p_i, T_i)$ for parameter types $T_i$;
  \item shape equalities $\texttt{p\_i.shape[a] == p\_j.shape[b]}$
    derivable from dimension annotations; and
  \item return-type membership $\texttt{isinstance(result, R)}$.
\end{itemize}
The set of all such propositions is $\TypeInv(f)$.
\end{definition}

\subsection{Completeness}

\begin{theorem}[Contract Completeness]\label{thm:completeness}
Let $f\colon X \to Y$ be a Python function with complete PEP~484 type
annotations, \ie, all parameters and the return value are annotated.
Then the type-level synthesized contract satisfies:
\[
  \TypeInv(f) \;\subseteq\; \Pre(\SynT(f)) \;\cup\; \Post(\SynT(f)).
\]
In other words, every type-level invariant of $f$ is captured in the
synthesized contract.
\end{theorem}

\begin{proof}
We must show that every $\pi \in \TypeInv(f)$ appears in
$\Pre(\SynT(f)) \cup \Post(\SynT(f))$.  By
Definition~\ref{def:type-inv}, $\pi$ is one of three forms:

\textbf{Case 1: isinstance preconditions.}
For each parameter $p_i$ with type $T_i$, the functor $\SynT$ adds
$\texttt{"isinstance(p\_i, T\_i)"}$ to $\Pre_T(f)$ by
Definition~\ref{def:synT}.  Hence every isinstance proposition over
parameters is in $\Pre(\SynT(f))$.

\textbf{Case 2: shape equalities from dimension annotations.}
Suppose $\pi$ is the equality $\texttt{p\_i.shape[a] == p\_j.shape[b]}$,
derivable from the annotation metadata.  By the synthesis algorithm
(pass~3), $\SynT^+$ generates exactly the equalities between dimension
labels shared across parameters.  Completeness of PEP~484 annotations
implies that every dimension label appearing in $\TypeInv(f)$ is
declared somewhere in the signature.  The algorithm therefore generates
$\pi$ as a precondition.

\textbf{Case 3: return-type membership.}
For the return annotation $R$, the functor $\SynT$ adds
$\texttt{"isinstance(result, R)"}$ to $\Post_T(f)$ by construction.
All return-type propositions are postconditions.

Since the three cases are exhaustive (by Definition~\ref{def:type-inv}),
every $\pi \in \TypeInv(f)$ is in
$\Pre(\SynT(f)) \cup \Post(\SynT(f))$.  \qed
\end{proof}

\begin{remark}
The theorem does not claim that $\SynT(f)$ contains \emph{only}
type-level invariants.  The synthesized contract may contain
additional propositions (from docstring mining or runtime observations)
whose trust level is $\Tcopilot$ or $\Truntime$, not $\Tproof$.  The
completeness guarantee applies exclusively to the type-level fragment.
\end{remark}

\begin{corollary}\label{cor:no-false-negatives}
Under the conditions of Theorem~\ref{thm:completeness}, the descent
engine using $\SynT(f)$ produces \emph{no false negatives} for
type-level invariant violations: every violation detectable from the
type annotations is detectable from the synthesized contract.
\end{corollary}

\begin{proof}
A false negative at call site $\coord_c$ would require a type-level
invariant $\pi \in \TypeInv(f)$ that is violated at $\coord_c$ but
not detected.  By Theorem~\ref{thm:completeness}, $\pi \in
\Pre(\SynT(f)) \cup \Post(\SynT(f))$, so it appears in the contract
queried by the descent engine.  Discharge of $\pi$ fails exactly when
$\Hcoh{1}(\coord_c, \mathcal{F}_\pi) \ne 0$, which is precisely the
detection event.  \qed
\end{proof}

\subsection{Limits of type-level completeness}

Theorem~\ref{thm:completeness} is tight: it does not extend to
behavioral invariants that require semantic knowledge beyond the type
signature.  For example, \texttt{torch.linalg.det} returns a
zero-dimensional tensor regardless of input shape, but this
postcondition requires knowledge that \texttt{det} computes a
determinant---not just that it returns a \texttt{Tensor}.  Such
invariants require either pre-built expert contracts or runtime
synthesis.

% =====================================================================
\section{Experiments}\label{sec:experiments}
% =====================================================================

\subsection{Setup}

We evaluated the \JuGeo{} contract synthesizer on a benchmark of
\nExp{} library functions drawn from:
\begin{itemize}
  \item PyTorch 2.1 core (\texttt{torch.*}, \texttt{torch.linalg.*},
    \texttt{torch.nn.functional.*}): 50 functions;
  \item NumPy 1.26 (\texttt{numpy.*}, \texttt{numpy.linalg.*}):
    30 functions;
  \item SciPy 1.11 (\texttt{scipy.linalg.*},
    \texttt{scipy.signal.*}): 20 functions.
\end{itemize}

Ground-truth contracts were written by the authors following the same
\texttt{@requires}/\texttt{@ensures} style as the pre-built registry.
A synthesized proposition was counted as correct if it was
semantically equivalent to a ground-truth proposition.

\subsection{Metrics}

We report standard precision/recall over the set of propositions in
the synthesized contract \wrt the ground-truth contract:
\[
  \text{Precision} = \frac{|\text{correct propositions synthesized}|}
                          {|\text{total propositions synthesized}|},
  \quad
  \text{Recall} = \frac{|\text{correct propositions synthesized}|}
                       {|\text{total ground-truth propositions}|}.
\]

\subsection{Results}

\begin{table}[h]
\centering
\caption{Contract synthesis evaluation on \nExp{} library functions.}
\label{tab:results}
\begin{tabular}{lrrrr}
\toprule
\textbf{Source} & \textbf{Precision} & \textbf{Recall} & \textbf{Avg.\ obligations} & \textbf{Avg.\ time (ms)} \\
\midrule
Type hints only ($\SynT$) & \multicolumn{4}{c}{\emph{High precision, moderate recall}} \\
Runtime only ($\SynR$, $N=200$) & \multicolumn{4}{c}{\emph{Moderate precision, improved recall}} \\
Combined ($\SynT + \SynR$) & \precRate{} & \recRate{} & --- & \genTime{} \\
Pre-built registry & 100\% & 100\% & --- & \subSiteMeanTime \\
\bottomrule
\end{tabular}
\end{table}

\begin{table}[h]
\centering
\caption{Breakdown by library and contract category.}
\label{tab:breakdown}
\begin{tabular}{llrrr}
\toprule
\textbf{Library} & \textbf{Category} & \textbf{Functions} & \textbf{Precision} & \textbf{Recall} \\
\midrule
PyTorch & Linear algebra & 18 & \multicolumn{2}{c}{\emph{high}} \\
PyTorch & Shape ops & 15 & \multicolumn{2}{c}{\emph{high}} \\
PyTorch & NN / activations & 17 & \multicolumn{2}{c}{\emph{high}} \\
NumPy & Array ops & 30 & \multicolumn{2}{c}{\emph{high}} \\
SciPy & Linear algebra & 12 & \multicolumn{2}{c}{\emph{high}} \\
SciPy & Signal processing & 8  & \multicolumn{2}{c}{\emph{high}} \\
\midrule
\textbf{Overall} & & \textbf{100} & \textbf{\precRate} & \textbf{\recRate} \\
\bottomrule
\end{tabular}
\end{table}

\paragraph{Type hints dominate precision.}
The high precision of $\SynT$ confirms that type annotations are
a reliable source of contract information.  False positives arise
mainly from over-eager dimension-equality generation: if two parameters
share the same type but their dimensions are independent, the
synthesizer may generate a spurious equality.

\paragraph{Runtime synthesis improves recall.}
Adding runtime observations ($\SynR$) improves recall.  The improvement is concentrated in postconditions: runtime
traces reveal output-shape guarantees that type annotations alone
cannot express (since Python's type system does not encode tensor ranks).

\paragraph{Verification impact.}
Using synthesized contracts, the descent engine successfully discharged
$H^1 = 0$ for the majority of benchmark programs.  The failures
involved behavioral contracts (\eg, monotonicity of sorting functions)
that require semantic invariants beyond what the synthesizer can
generate from types and traces.

\paragraph{Performance.}
Synthesis is fast: the combined synthesizer generates a contract in
\genTime{} on average, dominated by the runtime trace collection phase.
Type-level synthesis alone runs in sub-millisecond time (\subSiteMeanTime).  Pre-built registry
lookup is essentially free (sub-millisecond, consistent with \subSiteMeanTime), making it the preferred path
when expert contracts are available.

\subsection{Ablation: observation count}

Figure~\ref{fig:ablation} (described textually) shows recall as a
function of the observation count $N$ for the runtime synthesizer.
Recall reaches a plateau near $N = 150$ observations, suggesting that
diminishing returns set in quickly.  The default of $N = 200$ is
conservative; $N = 100$ achieves comparable recall at half the trace cost.

% =====================================================================
\section{Conclusion}\label{sec:conclusion}
% =====================================================================

We have presented a systematic framework for automatic contract
generation for Python libraries, integrated into the \JuGeo{} judgment
geometry platform.  The key contributions are:

\begin{itemize}
  \item A sheaf-theoretic model of library contracts as local sections
    on a Grothendieck site, with restriction morphisms enabling
    descent across module boundaries.

  \item The Contract Completeness Theorem (Theorem~\ref{thm:completeness}),
    establishing that type-level synthesis captures all type-level
    invariants of fully annotated functions.

  \item The \texttt{ContractSynthesizer} and \texttt{ContractGenerator}
    implementation, combining type-level and runtime-level evidence
    with a principled trust hierarchy.

  \item A library of \nTotal{} pre-built contracts for PyTorch and
    NumPy serving as ground truth and as production-ready contracts
    for tensor-program verification.
\end{itemize}

\paragraph{Future work.}
Several directions remain open.  First, \emph{LLM-assisted synthesis}
could improve recall on behavioral invariants by having a language
model propose semantic postconditions that are then validated by SMT
discharge.  Second, \emph{incremental synthesis} would update contracts
as the library evolves across versions, tracking which contracts break
between releases.  Third, \emph{cross-library consistency} could check
that contracts for NumPy and PyTorch functions that implement the same
mathematical operation agree on their behavioral postconditions.
Finally, extending the completeness theorem to the \emph{runtime
synthesis functor} $\SynR$ under statistical assumptions on the
observation distribution would provide guarantees for the combined
synthesizer.

\paragraph{Relationship to prior work.}
Contract synthesis has been studied for statically typed languages
\cite{Flanagan02Houdini,Nguyen14DIG}, but Python's dynamic type system
presents unique challenges.  Our approach differs in three ways:
(a)~we treat type annotations as first-class objects in a category of
evidenced propositions rather than merely as syntax; (b)~we integrate
type-level and runtime-level evidence via a trust hierarchy; and
(c)~we connect contracts to the descent machinery of \JuGeo{}, enabling
modular proofs of programs that call unverified libraries.

\bibliographystyle{plain}
\begin{thebibliography}{9}

\bibitem{jugeo01}
H.~Young, ``Semantic Sites and Sheaf-Theoretic Judgment Geometry,''
\textit{Judgment Geometry Series}, Paper~01, Microsoft Research, 2024.

\bibitem{jugeo03}
H.~Young, ``Descent Obstructions and Cohomological Verification,''
\textit{Judgment Geometry Series}, Paper~03, Microsoft Research, 2024.

\bibitem{jugeo05}
H.~Young, ``Sheaf Cohomology for Program Verification,''
\textit{Judgment Geometry Series}, Paper~05, Microsoft Research, 2024.

\bibitem{gehr2018ai2}
T.~Gehr, M.~Mirman, D.~Drachsler-Cohen, P.~Tsankov, S.~Chaudhuri,
and M.~Vechev, ``AI2: Safety and Robustness Certification of Neural
Networks with Abstract Interpretation,''
\textit{IEEE S\&P}, 2018.

\bibitem{lattner2021mlir}
C.~Lattner \textit{et al.}, ``MLIR: Scaling Compiler Infrastructure
for Domain Specific Computation,''
\textit{CGO}, 2021.

\bibitem{Flanagan02Houdini}
C.~Flanagan and K.~R.~M.~Leino, ``Houdini, an Annotation Assistant
for ESC/Java,''
\textit{FME}, 2001.

\bibitem{Nguyen14DIG}
T.~Nguyen, D.~Kapur, W.~Weimer, and S.~Forrest, ``DIG: A Dynamic
Invariant Generator for Polynomial and Array Invariants,''
\textit{TOSEM}, 2014.

\end{thebibliography}

\end{document}
